\section{Discussion}\label{sec:disc}

\subsection{Relation to earlier work}\label{ss:disc-prior}
For the decomposition of~\cite[\S4]{WuebbenII}, obtained from Hedden, Herald and Kirk's decomposition of $T(3,5)$
by regluing, the pillowcase Floer homology with zero bounding cochain already agrees with $\Inat(K_q)$ in every
grading, for odd $q\ge7$~\cite[Thm.~1.2]{WuebbenII} and for $q=5$ by the computation of Hedden, Herald and
Kirk~\cite[\S11]{HHK2}; no bounding cochain is needed there.

Smith proves that for $q=5$ his decomposition $\Qh{-1/2}\cup(Q_{1/3}+Q_{1/5})$ with zero bounding
cochains gives a Floer homology of rank $5$ rather than $\dim\Inat=7$ \cite[proof of Thm.~5.9]{Smith}, so that
under Conjecture~D of~\cite{CHKK} a nonzero bounding cochain is required \cite[Thm.~5.9]{Smith};
in his closing paragraph he proposes finding the bounding cochains that correct the rank.  Sections~\ref{sec:smith} and~\ref{sec:khovanov} give
two answers at $q=5$ and $q=7$ that, granting the arc case of the intersection formula~\cite[Thm.~5.25]{KWZ},
no rational closure distinguishes: the smoothings $(N_q,b)$, and their deformations
$(N_q,b+\kappa_q)\cong\mathrm{BN}_q$.

For rational tangles the traceless character varieties coincide with the invariants of Kotelskiy,
Watson and Zibrowius \cite[\S1.8]{KWZ}.  Hedden, Herald, Hogancamp and Kirk assemble Khovanov's tangle invariant from the
traceless character varieties of the crossingless resolutions \cite[Thm.~1.1 and Cor.~1.2]{HHHK},
with a result that agrees with $\widetilde{\Kh}(T)$ \cite{KWZagree}, and they ask whether higher terms
can be added so as to recover the Kronheimer--Mrowka spectral sequence \cite[Question~1.4]{HHHK}.
Cazassus, Herald, Kirk and Kotelskiy conjecture curve-valued instanton invariants $I(T)$ and
$I^\natural(T)$ of tangles, and predict that the relevant algebra is that of \cite{KWZ}
\cite[Conj.~E and \S11.4]{CHKK}.  They do not
conjecture $I(T)=\widetilde{\mathrm{BN}}(T)$, which cannot hold for all tangles because
$\dim\Inat<\rk\widetilde{\Kh}$ for some knots, for instance $T(4,5)$ \cite[\S\S10.2--10.3]{CHKK}.
Question~\ref{q:arc} concerns one tangle, $T_q$, for whose closures $P(-2,3,q)$ and $K_*$ the two
dimensions agree (\cite[\S3]{WuebbenII} and Proposition~\ref{prop:aux-closure}).  Kotelskiy, Watson
and Zibrowius suggest that the differentials of the spectral sequences from Khovanov homology to
Heegaard Floer and instanton homology might be realized by smoothing self-intersections of
$\widetilde{\mathrm{BN}}(T)$ and $\widetilde{\Kh}(T)$ \cite[\S1.8]{KWZ}.  Sections~\ref{sec:smith}
and~\ref{sec:khovanov} proceed in the opposite direction: a smoothing of the curve on the instanton
side reaches the Khovanov rank, and one further term carries it to the Bar-Natan invariant $\mathrm{BN}_q$.

Smoothing a self-intersection of an immersed curve by a degree-one bounding cochain is Lagrangian
surgery.  Palmer and Woodward prove that immersed Floer cohomology is invariant under surgery for
Lagrangians of dimension at least three, and in dimension two under an additional assumption
\cite[Thm.~1.3]{PalmerWoodward}; for curves on surfaces, Auroux and Smith compare surgery with the
mapping cone, which involves a bounding cochain \cite[Lem.~2.25 and Rem.~2.26]{AurouxSmith}.  The
corner term $\kappa_q$ is of a different kind: it is supported on chords around a corner
(Remark~\ref{rem:corner-term}), which are morphisms of the wrapped category of
\cite[\S11.5]{CHKK} and not self-intersection points of the curve.

\subsection{Open problems}\label{ss:disc-open}
\emph{The assigned object.}  Granting the arc case of the intersection formula of Kotelskiy, Watson and
Zibrowius~\cite[Thm.~5.25]{KWZ}, Question~\ref{q:arc} cannot be decided by rational closures at $q=5,7$
(Corollary~\ref{cor:closures-full}); curves parallel to $c$, and gradings, separate the
candidates (Remark~\ref{rem:sharpness}).  For $q\ge11$ a deformation of Smith's curve to
$\mathrm{BN}_q$, if it exists, is not small in the sense of Remark~\ref{rem:q11}.

\emph{The support of the bounding cochain.}  Proposition~\ref{prop:hyp-conflict} shows that
Hypothesis~\ref{hyp:q7-support} is not implied by closure data.  A direct proof would have to identify
the rigid figure-eight boundary strata of the quilted moduli spaces, exclude simultaneous support at
several self-intersection orbits, and intertwine the output with the switches of
\S\ref{ss:smoothings}.  In the strip-shrinking model the zero-width boundary stratum is predicted to
contain quilted trajectories with any positive number of figure-eight bubbles, and the zero-input
correction a count of figure-eight bubble trees \cite[\S\S4.2--4.4]{BottmanWehrheim}, so neither zero
input nor index zero forces support at a single self-intersection orbit.  A proof must use the geometry near the six internal
edge circles.

\emph{The formal route.}  Remark~\ref{rem:yoneda-recognition} reduces the identification of the
assigned object to a wrapped representability theorem of the kind Gao proves
\cite[Thm.~7.24 and Cor.~7.25]{Gao}, and Gao's theorem requires an embedded correspondence; by Proposition~\ref{prop:partial-triple} the correspondence
induced by the line of $Q_{1/3}$ has a triple point.  An immersed version of Gao's theorem covering
this triple point, or an embedded replacement inducing the same quilted module, would be needed, and
neither is available.  Independently, every selection statement of Sections~\ref{sec:smith}
and~\ref{sec:khovanov} is conditional on the correspondence between instanton and pillowcase
invariants \cite[Conj.~D]{CHKK}; neither Fukaya's nor Gao's correspondence theorem
\cite{FukayaCorr,Gao} compares the two.

\emph{Gradings of bounding cochains.}  In the normalization of Section~\ref{sec:khovanov} the object
$(N,b_{S_{69}})$ is bigraded, and the class of $b_{S_{69}}$ spans the quantum-degree-zero part of the
$\delta$-degree $-1$ endomorphism cohomology of $N$.  The $S_{18}/S_{25}$ and $S_{74}$ objects are not
bigraded (\S\ref{ss:smoothings}, Remark~\ref{rem:q7-scope}).

A selection by homogeneity of degree zero for the grading of $N$ would therefore choose the $S_{69}$ class, which the second closure excludes (Proposition~\ref{prop:aux-closure}).  Bigradedness of the deformed object alone does not select: among the objects allowed by Hypothesis~\ref{hyp:q7-support} with pairing dimension nine with $\mathcal E$, the bigraded ones also include $(N,b_{S_{18}}+b_{S_{74}})$ and the switches at $S_{34}$, $S_{38}$, $S_{75}$ and $S_{77}$, all with pairing dimension $25$ with $\mathcal E_*$.  None of the bigraded allowed objects has the pair $(9,31)$.  For $t>0$, with $(0,\pi)$ special (Remark~\ref{rem:sign-positive}), the same holds with $S_{10}$ in place of $S_{69}$.

No quantum grading is known on the instanton side, so this is an observation and not a test.  A
$\Z/4$-graded selection principle compatible with the absolute grading of Hedden, Herald and
Kirk~\cite[\S6.1]{HHK2} would replace the closure arguments of Section~\ref{sec:smith} by a structural one.
